\section{The cost of sharing a label on a Euclidean flat}
\label{sec:flatprice}
A \emph{spectral flat} is a Euclidean space equipped with spectral
prices, or an isometric copy of such a space inside a larger spectral
space. On a flat, the distances between pairs of sites determine the
price of a bit shared by \emph{all} the sites. The price is the mean
width of their convex hull (Theorem~\ref{thm:flatlaw}). For three sites
it is half the perimeter of their triangle, as in Figure~\ref{fig:triangle}.
We first show how to rescale prices on a flat and how to bound the cost
of passing an existing label to one more site.

\begin{lemma}[Homogeneous flat prices]\label{lem:homogenization}
Every spectral flat admits a compatible price system with the same
wire distances in which the price of each nonzero tensor is invariant
under translation, jointly convex in its leg positions, and positively
homogeneous under dilation.
\end{lemma}
\begin{proof}
Embed the flat in the geodesic spectral extension of
Theorem~\ref{thm:geometry}. Straight segments in the flat are geodesics,
so Corollary~\ref{cor:spectral-geodesic} gives joint convexity of
$H(T;x)$ in the leg positions $x$. Under a common translation along a
line, this is a convex function on $\mathbb R$. Transport from the first
site to each of the others bounds it above by a quantity depending only
on their relative positions, so it is constant. Next, $H(T;0)=0$,
and hence $H(T;tx)/t$ is nondecreasing for $t>0$. The same transport
bound gives a finite upper bound for this quotient. The limit as
$t\to\infty$ is therefore defined for every $T\ne0$ and every $x$.
The limiting prices inherit joint convexity and translation invariance,
and a change of scale gives positive homogeneity.
On each finite support of positive diameter, the rescaled price of its
binary diagonal lies between the diameter and a spanning-tree length.
Its limit is positive, and normalizing by it gives a spectral point
by compactness. Tensor prices on overlaps remain compatible. Wire
prices are the original Euclidean lengths at every scale.
\end{proof}

\subsection{The incremental cost of a common label}

Recall that $\delta(A)$ is the price of the binary diagonal on a finite
nonempty set of sites $A$. For a diagonal, the cost of adding one
site is exactly its increment in $\delta$. For a general tensor
carrying the same label at those sites, the increment gives an upper
bound. The following construction is related to the copying argument
with quadratic weights in \cite[Proposition~3.1]{nine}.

\begin{proposition}[Copying a common label]\label{lem:submod}
Let $I$ be a finite nonempty set of sites in a spectral space.
Suppose that $T$ carries the same $q$-valued basis label at each site
of $I$, so its coefficients vanish unless the marked indices agree.
Let $T'$ be the tensor obtained by appending a copy of this label at
$z\notin I$. Then
\[
 H(T')\le H(T)+(\delta(I\cup\{z\})-\delta(I))\log q.
\]
In particular $\delta$ is submodular on intersecting sets, that is,
$\delta(A\cup B)+\delta(A\cap B)\leq\delta(A)+\delta(B)$ if $A,B$ are finite and $A\cap B\ne\emptyset$.
\end{proposition}

\begin{proof}
If $I$ is a singleton, a wire from $I$ to $z$ suffices to append the
common label at $z$, and $\delta(I)=0$.

Write $I=\{1,\ldots,r\}$ with $r\geq2$. By
Lemma~\ref{lem:comparison} we may extend scalars to an algebraically
closed field while preserving the prices on the original tensors.
Write $h$ for the common label of $T$. Starting from
$T\otimes\Delta_{I\cup\{z\}}(L)$, we construct
$T'\otimes\Delta_I(q)$ by restriction and degeneration with
$L=O_r(q)$. The second factor in the output accounts for the
subtraction of $\delta(I)$ in the price bound.

Retain all indices of $T$, including $h$ at each site in $I$, and
append an index $v_j$ at site $j\in I$ and an index $g$ at $z$. All
these labels range from $0$ to $q-1$. We first impose the
single integer equation
\[
  g=h+\sum_{j=2}^r v_j-(r-1)v_1
\]
by coefficient extraction. In an indeterminate $t$, assign a monomial
to each choice of local labels. Its exponent is $q-1-h+(r-1)v_1$ at
site~1, $q-1-v_j$ at site $j$ for $2\le j\le r$, and $g$ at $z$.
The coefficient of $t^{r(q-1)}$ in their product is the indicator of
the equation. The product has degree at most $2r(q-1)$, so interpolation
expresses this coefficient as a linear combination of
$L=2r(q-1)+1$ evaluations. Use the common label of the supplied
$L$-diagonal to choose the same evaluation point at all these sites,
with the interpolation weight included at site~1. This is a restriction.
The exponents are integers in every characteristic.

On this support, $(g-h)^2$ is a sum of local integer weights. Indeed
in $g^2-2hg+h^2$ one replaces $g$ in the cross term by the right side
of the equation, using the copy of $h$ at each site of $I$.
Multiply the local basis vectors by the corresponding powers of a
nonzero parameter. The local maps may have Laurent entries, but every tensor
coordinate extends polynomially to zero because every total exponent on the
support is nonnegative. Specializing to zero retains exactly $g=h$. The
remaining equation is $\sum_{j=2}^r v_j=(r-1)v_1$. A second such
specialization uses the local weight sum
\[
  \sum_{j=2}^r v_j^2-(r-1)v_1^2=\sum_{j=2}^r(v_j-v_1)^2
\]
on the remaining support. It retains exactly the terms in which all
$v_j$ agree. Their common value is independent of $h$, giving the
claimed product.
Lemma~\ref{lem:spectral-closure} applied to the two specializations
gives
\[
 H(T')+\delta(I)\log q\le H(T)+\delta(I\cup\{z\})\log L .
\]
Apply the construction to $T^{\otimes m}$ with $q^m$ labels.
Since $m^{-1}\log(2r(q^m-1)+1)\to\log q$, division by $m$ gives the
copying inequality. For the last assertion, apply this inequality to
the diagonal on $A$ with marked sites $I\subseteq A$.
Starting with $I=A\cap B$ and appending the sites of $B\setminus A$
to both $I$ and $A$ then gives the proposition by telescoping.
\end{proof}


For two existing copies at $a,b$, the price increment per bit is at
most half the extra distance incurred by going through $z$,
\[
 \tfrac12\bigl(d(a,z)+d(z,b)-d(a,b)\bigr).
\]
Indeed, the three identity wires of dimension $n$ around the triangle
supply a diagonal of size $n^{2-o(1)}$ by the following restriction
\cite{Strassen87,Behrend,cvzgraph}. Let the three edge labels be
$i,j,k\in[n]$, and take a set $B\subseteq[\lfloor n/3\rfloor]$ of
size $n^{1-o(1)}$ with no nonconstant three-term arithmetic progression.
At the three sites retain coordinate pairs with
$j-i\in B$, $k-j\in B$, and $k-i\in \{2u:u\in B\}$, respectively.
The two consecutive differences must then be equal.
Each surviving triple is determined by any one of its pairs. There are
at least $\lfloor n/3\rfloor |B|$ such triples, by taking
$i\le\lfloor n/3\rfloor$ and the common difference in $B$. Taking
prices and dividing by $\log n$ gives
$2\delta(\{a,b,z\})\le d(a,b)+d(b,z)+d(z,a)$. The claimed increment
follows from Proposition~\ref{lem:submod} and $\delta(\{a,b\})=d(a,b)$.

\subsection{Mean width on Euclidean flats}

For a nonempty polytope $K\subset\mathbb R^D$, $D\ge1$, let
$s_K(u)=\max_{x\in K}\langle u,x\rangle$ be its support function.
The length of the projection of $K$ onto a unit vector $u$ is
$s_K(u)+s_K(-u)$. The \emph{mean width} averages these lengths over
directions, with a constant chosen so that a segment has its own length
as mean width. Equivalently,
\[
 W(K)=\sqrt{2\pi}\,\mathbb E\,s_K(G)
     =\int_{\mathbb S^{D-1}}s_K(u)\,d\sigma_D(u),
 \qquad G\sim N(0,I_D).
\]
Here $\sigma_D$ is the rotation-invariant measure with this
normalization. The two expressions agree by radial symmetry and
$\sqrt{2\pi}\,\mathbb E(G_1)_+=1$, where
$a_+=\max(a,0)$. The Gaussian expression is unchanged by translation
and by isometric embedding into a higher-dimensional Euclidean space.
A singleton has width zero. A finite nonempty set and its convex hull
have the same support function and hence the same mean width.

\begin{theorem}[Mean width of a shared label]\label{thm:flatlaw}
Suppose $\mathbb R^D$ embeds isometrically in a spectral space. For
every nonempty finite set $V\subseteq\mathbb R^D$, the price of its
binary diagonal is $\delta(V)=W(V)$. Consequently the price of a
common $q$-valued label on $V$ is $(\log q)W(V)$.
\end{theorem}

\begin{proof}
For fixed $V$, the price of the diagonal with one additional leg is a
convex function of that leg's position. At each site of $V$ the new
leg can be copied locally, so all these values equal $\delta(V)$.
Throughout their convex hull the price is therefore at most $\delta(V)$,
and discarding the added leg gives the reverse inequality. We may thus
write $\delta(K)$ for the price of any finite set with convex hull $K$.

\emph{Minkowski sums.}
First assume $\delta(tK)=t\delta(K)$ for $t>0$. The translation argument
in Lemma~\ref{lem:homogenization} applies to the diagonal. To see
subadditivity for Minkowski addition, list the vertices $(x_i)$ of $K$
and $(y_j)$ of $L$. Take a diagonal with one leg for each pair $(i,j)$.
Placing this leg at $x_i$ gives price $\delta(K)$, and placing it at
$y_j$ gives price $\delta(L)$. At the midpoint placement the convex
hull is $\tfrac12(K+L)$. Joint convexity and homogeneity give
$\delta(K+L)\leq\delta(K)+\delta(L)$. In particular, $\delta$ is
Lipschitz in Hausdorff distance. Polytopes within distance
$\varepsilon$ of each other are each contained in the sum of the other
with $\varepsilon$ times a fixed cube containing the unit ball.

\emph{Segments and triangles.}
Adding a segment increases the price by exactly its length. For the
vertex set $S$ of $K$, form $\bigcup_{j=0}^m(S+jv)$ by adjoining
$(S+(j-1)v)\cup(S+jv)$ one step at a time. At step $j$ the overlap
contains $S+(j-1)v$. Proposition~\ref{lem:submod} and translation invariance
give
\[
 m|v|\leq\delta(K+[0,mv])
 \leq \delta(K)+m\big(\delta(K+[0,v])-\delta(K)\big).
\]
The first inequality uses a pair of points distance $m|v|$ apart.
Let $m\to\infty$ and use subadditivity to conclude
$\delta(K+[0,v])=\delta(K)+|v|$.

Next let $T$ be a nondegenerate triangle. The wire restriction above
gives $\delta(T)\le\tfrac12\operatorname{per}(T)=W(T)$.
For the last equality, write $e_1,e_2,e_3$ for the oriented edge
vectors. The length of the projection onto a unit direction is half
the sum of the absolute edge projections. By symmetry of $G$,
\[
 W(T)=\frac{\sqrt{2\pi}}{4}\sum_{i=1}^3\mathbb E|\langle G,e_i\rangle|
      =\frac12\sum_{i=1}^3\|e_i\|.
\]
Moreover $T+(-T)$ is the sum of
the three segments given by those oriented edges. Indeed, the sums of
subsets of the three edge vectors give zero and all pairwise vertex
differences. Segment additivity gives
\[
\delta(K)+2W(T)=\delta(K+T+(-T))
       \le \delta(K+T)+W(T)\le \delta(K)+2W(T).
\]
Here the inequalities use subadditivity and the half-perimeter bound
for both triangles. Equality throughout gives
$\delta(K+T)=\delta(K)+W(T)$. The same identity for degenerate
triangles follows from segment additivity and translation invariance.

\emph{General polytopes.}
By Lemma~\ref{lem:triangle-density}, we can approximate $s_K$
uniformly by a linear combination of triangle support functions.
Grouping the positive and negative coefficients expresses the combination
as $s_{P_+}-s_{P_-}$ for two Minkowski sums of scaled triangles. If the
error is at most
$\varepsilon$, then $K+P_-$ and $P_+$ are within Hausdorff distance
$\varepsilon$. Using the Lipschitz bound for $\delta$, its additivity
along triangles, and the additivity and continuity of $W$ gives
$\delta(K)=W(K)$ as $\varepsilon\to0$.

For the original price, $\delta(tK)/t$ is nondecreasing in $t$.
The proof of Lemma~\ref{lem:homogenization} produces a homogeneous
spectral price system at either endpoint, $t\downarrow0$ or
$t\to\infty$. Indeed, for a finite set of positive diameter, its
rescaled diagonal price has the same positive lower bound in both
limits. The homogeneous case identifies both endpoint values as
$W(K)$, so the quotient has this value at $t=1$. The diagonal scaling
rule of Section~\ref{sec:spectra} gives the factor $\log q$ for the
$q$-valued label.
\end{proof}

The density statement used above also follows from Schneider's triangle
theorem \cite{Schneider1996}, as formulated by Schneider and Schuster
\cite[pp.~2--3]{SS}. We give a proof by testing against signed measures
on the sphere.
\begin{lemma}[Approximation by triangle support functions]
\label{lem:triangle-density}
The support functions of triangles in $\mathbb R^D$, including degenerate
triangles, span a uniformly dense subspace of $C(\mathbb S^{D-1})$.
\end{lemma}

\begin{proof}
Suppose a finite signed measure $\mu$ on the sphere annihilates all
triangle support functions. For $v,w\in\mathbb R^D$ and $h>0$ it
annihilates the difference quotient
\[
 \frac{\max(0,\langle v,u\rangle,h\langle w,u\rangle)
              -\max(0,\langle v,u\rangle)}{h}.
\]
As $h\downarrow0$, the quotient converges boundedly to
$\mathbf1_{\{\langle v,u\rangle\leq0\}}(\langle w,u\rangle)_+$.
Using both $w$ and $-w$ shows that the restriction $\nu$ of $\mu$
to any closed hemisphere annihilates $|\langle w,u\rangle|$ for
every $w$.

To determine the even part of $\nu$, average in $w=G+tz$ for a standard
Gaussian $G$ and a fixed vector $z$. With $b(a)=\mathbb E|G_1+a|$ this
gives $\int b(t\langle z,u\rangle)\,d\nu(u)=0$. Since $b(0)>0$ and
$b''(a)=\sqrt{2/\pi}\exp(-a^2/2)$, every even derivative at zero is
nonzero. Differentiating under the integral gives
$\int\langle z,u\rangle^{2j}\,d\nu(u)=0$ for all $j\ge0$.
Polarization and polynomial approximation give vanishing against all
continuous even functions. Hence $\nu$ is odd under $u\mapsto-u$.
It is supported on a closed hemisphere, so it vanishes in the interior
of that hemisphere. Such interiors cover the sphere, giving $\mu=0$.
The Riesz representation and Hahn--Banach theorems prove the density.
\end{proof}


For any planar convex polygon $K$, the two boundary chains between its
extreme projected points each project monotonically. Thus its projected
length is half the sum of the absolute edge projections, just as for a
triangle. Averaging gives $W(K)=\tfrac12\operatorname{per}(K)$, with
a segment counted twice in the perimeter \cite{Schneider}.

Theorem~\ref{thm:flatlaw} also evaluates the increment in
Proposition~\ref{lem:submod}. If a tensor on a flat already carries a
common label at $I$, copying it to any $z\in\operatorname{conv}(I)$
preserves the price. The convex hull and hence its mean width stay the
same, and discarding the new copy recovers the original tensor.
